\documentclass[11pt,letterpaper]{article}
\usepackage[T1]{fontenc}
\usepackage{lmodern,amsmath,amssymb,amsthm,microtype}
\usepackage[margin=1in]{geometry}
\usepackage[hidelinks]{hyperref}
\newtheorem{theorem}{Theorem}
\title{The minimum size of a maximal partial arrangement of Euler's 36 officers}
\author{Aafi Khan Malek\\
\small Independent researcher\\
\small\href{mailto:aafimalek2032@gmail.com}{\texttt{aafimalek2032@gmail.com}}}
\date{27 November 2026}
\begin{document}
\maketitle

\begin{abstract}
The smallest maximal partial arrangement of Euler's 36 officers has 18
officers. The upper construction was already known. We prove the lower bound
by a finite case split and independently replayed SAT certificates. In
particular, the size-17 exclusion uses 56 normalized instances. The
certificate generators and case reductions accompany the paper. The upper
construction is verified in Lean; the full lower bound is not.
\end{abstract}

\section{Problem and result}
Each officer has one of six ranks and one of six regiments, with each
rank--regiment pair available once. A partial arrangement in a $6\times6$
grid repeats neither a rank nor a regiment within any row or column. It is
\emph{maximal} if no further officer can be added without breaking a rule.
This is a different question from the classical problem of placing all 36.
Equivalently, we seek the least number of filled cells in a maximal pair of
orthogonal partial Latin squares of order six \cite{dgy25}.

\begin{theorem}
Every maximal partial arrangement contains at least 18 officers, and one
contains exactly 18.
\end{theorem}

\paragraph{The known upper bound.}
For $c=0,1$ and $a,b\in\mathbf F_3$, use the words
\[
 (3c+a,\,3c+b,\,3c+(a+b\bmod 3),\,3c+(a+2b\bmod 3)).
\]
These words encode row, column, rank, and regiment. In either block any two
of the four linear forms are independent, so every pair of labels in
different coordinates of that block occurs exactly once. An additional
word would contain two labels in the same block, so it repeats a pair.
This is the $3+3$ upper construction recorded by Donovan, Grannell and
Yaz\i c\i{} \cite[Table~1]{dgy25}. Their reported minimality checks were
for orders four and five. The later general exact results start at order
21 \cite{dgy26,dgy26b}. Our contribution is the order-six lower bound.

\section{The lower bound}
We view placements as words in $\{0,\ldots,5\}^4$. A used pair is a pair of
labels in distinct coordinate groups that occurs in a selected word. A
valid design uses each pair at most once. It is maximal exactly when every
candidate word contains a used pair. Each coordinate group has six labels
and the degrees of its labels sum to the number of selected words.

The \href{https://github.com/Aafimalek/novel-discoveries/tree/research/euler-order16/euler-36-officers-disconnected}{companion
repository} gives the size-14 and size-15 exclusions and the size-16 audit.
The size-14 structural lower bound, before those exclusions, is separately
formalized in Lean. We treat the remaining size~17 below.

\paragraph{Unused labels.}
Suppose one group contains a degree-zero label. On the other three groups,
join labels whose pair is unused. The resulting tripartite graph is
triangle-free by maximality. A vertex of placement degree $d$ has the
same missing degree $e=6-d$ to either other group. Along any missing edge
$uv$, their neighbor sets in the third group are disjoint, so
$e(u)+e(v)\le6$. There are $M=36-17=19$ missing edges between any two groups.
Summing over these edges gives a combined squared-degree sum at most
$6M$. The two sums of degrees are both $M$, so Cauchy--Schwarz gives a
combined squared-degree sum at least $2M^2/6$, forcing $M\le18$.
This rules out degree zero. A corresponding lower bound follows in Lean
from the earlier matching inequality.

\paragraph{Degree one.}
With minimum degree one, choose a degree-one label and normalize its
unique word. If the other three labels in the word all have degree one,
it is an isolated component, impossible by the disconnected lower bound.
Otherwise classify by the three partner degrees, increasing. In the
companion audit, the four cases with maximum partner degree at least
three, the cases of types $112$ and $122$, and the normalized type
$222$ cases yield respectively $4$, $14$, $2$, and $7$ formulas.
All $27$ formulas were refuted.

\paragraph{Degree two.}
The remaining minimum is two, and every group contains a degree-two label
because its degree sum is 17. Call such a label low. For an unused low-low
pair, compare the two sets of used neighbors in each remaining group.
The following cases are exhaustive, in the given order:
\begin{enumerate}
\item All low-low pairs are used. There is one low in each group and
      the four lows share a word. Their four other words give $17$
      normalized diagrams.
\item Some unused low-low pair has two common neighbors in one group.
      The candidate rectangle bounds force no overlap in the other group,
      and there is one normalized formula.
\item Some unused low-low pair has a common neighbor in both remaining
      groups. There are three overlap diagrams. For profiles with at
      least four degree-at-least-three labels in each group these give
      three formulas. Otherwise a group has at least three lows;
      normalizing that group either to an anchor group or to a remaining
      group gives six further formulas.
\item Some unused low-low pair has just one common neighbor. This gives
      one formula.
\item Unused low-low pairs exist but none shares a neighbor in another
      group. This also gives one formula.
\end{enumerate}
These cases contribute $29$ refutations. The mathematical
normalizations, each specified instance, and their exact coverage appear
in the audit. A larger minimum degree would force at least 18 words
directly from one group's degree sum.

\section{Independent verification}
The CNF encoding selects exactly 17 words, defines pair use from those
words, enforces pair uniqueness, and prohibits any candidate with all
six pairs unused. The case constraints fix labels or enforce necessary
degree and rectangle conditions justified in the documentation. The
\href{https://github.com/Aafimalek/novel-discoveries/tree/research/euler-order16/euler-36-officers-disconnected/verification/order17-exact}{size-17
audit} includes a manifest with the exact DIMACS and proof hashes,
reproduction commands, and an independent verifier. All $56$
size-17 refutations passed DRAT-trim after generation with Kissat.
The certificate binaries are generated on demand, not bundled.
The upper witness was also checked in Lean 4.36.0 in
\texttt{EulerOfficers.exists\_maximal\_eighteen}. The lower-bound SAT
certificates and the full encoding have not been checked in Lean.
This note has not been peer reviewed.

\paragraph{Contributions and AI disclosure.}
Aafi Khan Malek selected and directed the research project and requested
the verification and publication workflow. ChatGPT assisted with this
extension's case analysis, computational proof, manuscript, and local
verification. Codex was not used for this extension. The companion
repository retains the attribution of earlier work.

\begin{thebibliography}{9}
\bibitem{dgy25} D. M. Donovan, M. J. Grannell and E. \c{S}. Yaz\i c\i,
On maximal orthogonal partial Latin squares and minimal codes with
specified length, minimum distance and covering radius,
\emph{Des. Codes Cryptogr.} (2025).
\href{https://doi.org/10.1007/s10623-025-01704-x}{DOI: 10.1007/s10623-025-01704-x}.
\bibitem{dgy26} D. M. Donovan, M. J. Grannell and E. \c{S}. Yaz\i c\i,
On the minimum number of entries in a pair of maximal orthogonal partial
Latin squares, \href{https://arxiv.org/abs/2602.09908}{arXiv:2602.09908} (2026).
\bibitem{dgy26b} D. M. Donovan, M. J. Grannell and E. \c{S}. Yaz\i c\i,
On the minimum number of entries in a set of maximal mutually orthogonal
partial Latin squares, \href{https://arxiv.org/abs/2611.13961}{arXiv:2611.13961} (2026).
\end{thebibliography}
\end{document}
